\section{Introduction}
\label{sec:intro}

A manipulation policy trained on one robot is expected to keep working when the robot changes. Hardware evolves faster than task-specific data can be collected, and every new gripper or arm restarts the collection effort when the policy has to be retrained from scratch. Vision-language-action models promise a way out by pretraining on many robots and transferring to new ones. Turning that promise into a measurable property requires an answer to a simple question: what does a policy carry over from the robots it was trained on, and which parts of that baggage stop it from working on a robot outside its training set.

This paper studies zero-shot cross-embodiment transfer as a controlled experiment around that question. Two obstacles have kept the literature from answering it. The first is protocol. Reports of zero-shot transfer fall into two groups that differ in one condition. In one group the target robot is absent from every stage of training. In the other group the target robot appears during pretraining and is only absent from task-specific post-training. The two settings answer different questions, and treating them as one attributes to a model design what a data composition produced. The second obstacle is confounding. Embodiment changes in existing evaluations arrive together with changes of task, scene, camera, data quantity and protocol, so a failure on a new robot has several possible causes. We therefore hold the manipulation regime fixed to stationary tabletop manipulation with fixed-base arms and two-finger grippers, and we vary the robot alone.

Within this regime we organise the study around four channels through which the source robots leave their mark on a policy. The state-action representation decides in which frame the policy reads its state and writes its action, and a frame that belongs to the source robot ties the policy to that robot. The diversity of the source family decides how much of that frame the policy can marginalise away. The auxiliary supervision decides which structure of the task the policy is pushed to represent beyond imitation. The exposure of the target during pretraining decides whether the policy has seen the target at all. Each channel becomes one research question, and every question is answered under the same backbone, the same pretraining budget and the same evaluation.

The representation channel receives the most attention because it admits a hierarchy. A motion of the gripper can be described relative to the robot base, relative to the gripper itself, or relative to the camera that observes it. The base is a convention of the manufacturer. The gripper frame is shared by every two-finger robot up to the geometry of its fingers. The camera is the one frame that every embodiment in a matched setup shares exactly, because the image is the channel through which the policy perceives the world. Describing the action in the gripper frame is a known aid to transfer. This study adds the last step of the hierarchy: the action is described by its effect, the motion that the fingertips and pads of the gripper undergo in the image over the next fraction of a second. The starting position of each keypoint follows analytically from the state of the robot and its kinematics, the policy predicts only the displacement, and a new robot decodes a command from the predicted displacement through its own kinematics. Embodiment enters at the input boundary through the anchors and at the output boundary through the decoding, and the learned mapping in between carries no convention of the source robot. The same quantity also serves as an auxiliary target for a policy that keeps emitting gripper-frame commands, which places it on the supervision channel as well.

To measure each channel we build a controlled benchmark in simulation. A procedurally generated family of 128 Franka-type arms with varied link lengths, finger shapes and colours is the source for pretraining under a fixed budget of 48K trajectories. Eight commercially used embodiments, all fixed-base arms with two-finger grippers, are held out from all training and grouped by what changes: appearance, gripper, arm, or both. Cameras are aligned in a shared tabletop frame so that image differences come from the robot itself. Post-training uses only the source Franka on three downstream tasks, and evaluation scores task progress on every target with thousands of rollouts per model. The whole study fits sixteen GPUs over three weeks, which keeps the design within reach of an academic group.

The four channels turn out to matter in a consistent order. Moving state and action from the base frame to the gripper frame raises average progress by eighteen points. Moving the action into the camera frame adds two more points on average and five on full embodiment changes when the command is decoded from the effect, and seven points on average and thirteen on full changes when the effect is an auxiliary target next to gripper-frame commands. Under a fixed budget, 128 source embodiments beat a single source by sixteen points, and effect supervision at eight sources matches gripper-frame commands at 128. Among auxiliary targets, those that ground the interaction in the image help most, and the effect target leads language, waypoint and object-box targets. Seeing the target during pretraining raises every representation and narrows the gaps between them, which confirms that the two protocols must be reported apart. The prediction error of the effect on a target orders the targets by their transfer, and the failures that effect supervision removes are the grasp-phase failures in which the geometry of the gripper decides the outcome.

Our contributions are the following. We separate target-absent from target-seen zero-shot transfer and show empirically that the two protocols measure different quantities. We present a controlled benchmark that isolates appearance, gripper, arm and full changes among fixed-base two-finger manipulators at an academic compute scale. We provide a factorised study of representation frames, source diversity, auxiliary targets and target exposure under one backbone and one budget. We introduce camera-frame motion effects as the most embodiment-free point of the representation hierarchy, together with a decoder that makes them executable on an unseen gripper and a diagnostic that anticipates transfer before any task rollout.
